% EDIT HERE: title, subtitle, author, institution and accent colour.
% Self-contained 16:9 slides. The example diagram uses no external images.
\documentclass[aspectratio=169,11pt]{beamer}
\usepackage[T1]{fontenc}
\usepackage{tgheros}
\usepackage{booktabs}
\definecolor{accent}{HTML}{245C61}
\definecolor{ink}{HTML}{182E3B}
\definecolor{muted}{HTML}{60717B}
\definecolor{wash}{HTML}{EDF4F3}
\setbeamercolor{normal text}{fg=ink,bg=white}
\setbeamercolor{structure}{fg=accent}
\setbeamercolor{frametitle}{fg=ink}
\setbeamercolor{title}{fg=ink}
\setbeamercolor{subtitle}{fg=accent}
\setbeamercolor{block title}{fg=white,bg=accent}
\setbeamercolor{block body}{fg=ink,bg=wash}
\setbeamercolor{alerted text}{fg=accent}
\setbeamerfont{title}{size=\Huge,series=\bfseries}
\setbeamerfont{subtitle}{size=\large}
\setbeamerfont{frametitle}{size=\LARGE,series=\bfseries}
\setbeamertemplate{navigation symbols}{}
\setbeamertemplate{itemize item}{\small\raise1pt\hbox{\textbullet}}
\setbeamertemplate{itemize subitem}{\textendash}
\setbeamersize{text margin left=11mm,text margin right=11mm}
\setbeamertemplate{frametitle}{\vspace{6mm}\insertframetitle\par\vspace{3mm}{\color{accent}\hrule height 0.7pt}}
\setbeamertemplate{footline}{%
\hspace{11mm}\color{muted}\scriptsize\insertshorttitle\hfill\insertframenumber\,/\,\inserttotalframenumber\hspace{11mm}\vspace{4mm}}
\title[Presentation title]{Presentation title}
\subtitle{One clear sentence about your contribution}
\author{Your name}
\institute{Your department / Institution}
\date{\today}

\begin{document}
\begin{frame}[plain]
\vspace{8mm}
{\small\color{accent}\insertinstitute}\par
\vspace{7mm}
{\usebeamerfont{title}\inserttitle\par}
\vspace{4mm}
{\usebeamerfont{subtitle}\color{accent}\insertsubtitle\par}
\vspace{7mm}{\color{accent}\hrule height 1.5pt}\vspace{5mm}
\insertauthor\par
{\small\color{muted}\insertdate}
\end{frame}

\section{The question}
\begin{frame}{Start with the question}
\begin{block}{Research question}
What changes when we redesign a research workflow?
\end{block}
\vspace{4mm}
\begin{itemize}\setlength{\itemsep}{5mm}
\item \textbf{Context:} Explain the problem and who experiences it.
\item \textbf{Gap:} State what existing approaches leave unresolved.
\item \textbf{Contribution:} Describe the new evidence or method you offer.
\end{itemize}
\end{frame}

\begin{frame}{Give the audience a route}
\begin{columns}[T,onlytextwidth]
\column{0.30\textwidth}
{\Huge\color{accent}01}\par\vspace{3mm}
\textbf{Question}\par\vspace{2mm}
What matters, and what do we need to learn?
\column{0.30\textwidth}
{\Huge\color{accent}02}\par\vspace{3mm}
\textbf{Evidence}\par\vspace{2mm}
How was the study designed, and what changed?
\column{0.30\textwidth}
{\Huge\color{accent}03}\par\vspace{3mm}
\textbf{Implications}\par\vspace{2mm}
What can we conclude, and what comes next?
\end{columns}
\end{frame}

\section{Approach}
\begin{frame}{Make the method easy to follow}
\begin{columns}[T,onlytextwidth]
\column{0.46\textwidth}
\textbf{Study design}\par\vspace{3mm}
\begin{enumerate}\setlength{\itemsep}{4mm}
\item Define a baseline and revised workflow.
\item Use the same task and outcome measures.
\item Compare time, errors and uncertainty.
\end{enumerate}
\column{0.46\textwidth}
\begin{block}{What to report}
Sample and selection criteria\par\medskip
Procedure and controls\par\medskip
Analysis and limitations
\end{block}
\end{columns}
\end{frame}

\section{Evidence}
\begin{frame}{Show the result, then explain it}
\begin{columns}[c,onlytextwidth]
\column{0.54\textwidth}
\centering
\begin{tabular}{@{}lrr@{}}
\toprule
& Baseline & Revised\\
\midrule
Time (min) & 24 & 18\\
Error rate (\%) & 8 & 5\\
Tasks & 40 & 40\\
\bottomrule
\end{tabular}
\par\vspace{3mm}{\scriptsize\color{muted}Illustrative values; replace with your own data.}
\column{0.38\textwidth}
{\Huge\bfseries\color{accent}25\%}\par\vspace{2mm}
less completion time in this example.\par\vspace{4mm}
Explain whether the difference is meaningful and how certain the estimate is.
\end{columns}
\end{frame}

\begin{frame}{Separate the finding from its limits}
\begin{columns}[T,onlytextwidth]
\column{0.46\textwidth}
\begin{block}{What the evidence supports}
State your strongest finding in one sentence. Connect it directly to the
research question and the evidence on the previous slide.
\end{block}
\column{0.46\textwidth}
\textbf{What remains uncertain}\par\vspace{3mm}
\begin{itemize}\setlength{\itemsep}{3mm}
\item Scope and sample size.
\item Measurement assumptions.
\item Generalization to other settings.
\end{itemize}
\end{columns}
\end{frame}

\section{Takeaways}
\begin{frame}{Leave three things to remember}
\begin{enumerate}\setlength{\itemsep}{6mm}
\item \textbf{The answer.} Respond directly to the research question.
\item \textbf{The contribution.} Say what your work makes possible.
\item \textbf{The next step.} Propose a concrete action or follow-up study.
\end{enumerate}
\end{frame}

\begin{frame}[plain]
\vfill
{\Huge\bfseries Questions?}\par\vspace{5mm}
{\large\color{accent}Thank you for your attention.}\par
\vspace{8mm}{\color{accent}\hrule height 1.5pt}\vspace{5mm}
\insertauthor\quad /\quad\href{mailto:you@example.com}{you@example.com}\par
{\small\color{muted}Add a link to your paper, code or supporting material.}
\vfill
\end{frame}
\end{document}
